\documentclass{report}
\usepackage{amsmath,amssymb}
\setlength{\parindent}{0mm}
\setlength{\parskip}{1em}
\begin{document}
\begin{center}
\rule{6in}{1pt} \
{\large Ulrich Ruede \\
{\bf A multi-scale algorithm to recover optimal convergence in the presence of corner singularities}}

Universitaet Erlangen \\ Cauerstrasse 11 \\ 91080 Erlangen \\ Germany
\\
{\tt ruede@cs.fau.de}\\
Christian Waluga\\
Barbara Wohlmuth\\
Herbert Egger\end{center}

Elliptic problems with reentrant corners exhibit the well-known
pollution-effect, i.e. the convergence deteriorates in global norms. This
can be avoided by many techniques, such as a suitable mesh grading or by
enriching the finite element spaces. Here we analyze a new method that
restores optimal converge in weighted norms and that differs form
alternatives in requiring only local corrections on uniform or
quasi-uniform grids. We show that the corrections must be such that the
modified finite element discretization reproduces the energy of the
solution with sufficient accuracy. This can be achieved by a suitable
modification in a small, fixed number of elements. The corresponding
stiffness matrix differs from the uncorrected one in only a few numerical
values. The correction parameter for each reentrant corner can be
computed efficiently in a multi-scale approach that can be �integrated in
a full multigrid method.


\end{document}
